\documentclass[12pt,a4paper,oneside]{report}

\usepackage[T1]{fontenc}
\usepackage[utf8]{inputenc}
\usepackage[french]{babel}
\usepackage[a4paper,margin=2.5cm]{geometry}
\usepackage{setspace}
\usepackage{graphicx}
\usepackage{float}
\usepackage{booktabs}
\usepackage{longtable}
\usepackage{array}
\usepackage{caption}
\usepackage{listings}
\usepackage{xcolor}
\usepackage[hidelinks]{hyperref}
\usepackage{enumitem}
\usepackage{fancyhdr}
\usepackage{tikz}
\usetikzlibrary{arrows.meta,positioning,shapes.geometric}

\onehalfspacing
\setcounter{secnumdepth}{3}
\setcounter{tocdepth}{3}

\pagestyle{fancy}
\fancyhf{}
\fancyhead[L]{\nouppercase{\leftmark}}
\fancyhead[R]{\thepage}

\definecolor{codegray}{gray}{0.95}
\lstdefinestyle{mystyle}{
  backgroundcolor=\color{codegray},
  basicstyle=\ttfamily\small,
  breaklines=true,
  frame=single,
  columns=fullflexible
}
\lstset{style=mystyle}

\begin{document}

% -------------------------
% Page de garde
% -------------------------
\begin{titlepage}
  \centering
  {\Large ANNEE UNIVERSITAIRE 2025 -- 2026\par}
  \vspace{1.5cm}
  {\Huge\bfseries RAPPORT DE PROJET DE FIN D'ETUDES\par}
  \vspace{1cm}
  {\Large IHEC CARTHAGE\par}
  \vspace{1.5cm}
  {\LARGE\bfseries Conception et developpement d'un assistant touristique intelligent multilingue pour les hotels tunisiens\par}
  \vspace{1.5cm}
  {\large Organisme d'accueil : SPECTRUM TECHNIQUE ET INFORMATIQUE\par}
  \vspace{1cm}
  {\large Elabore par :\par}
  {\large Chedly Hachicha\par}
  {\large Aziz Guibene\par}
  \vfill
  \begin{flushleft}
    \textbf{Encadrants :} Ines Hammami Hafsa et Youssef Bouzouaya
  \end{flushleft}
\end{titlepage}

\pagenumbering{roman}

\chapter*{Dedicaces}
\addcontentsline{toc}{chapter}{Dedicaces}
Texte de dedicaces a personnaliser.

\chapter*{Remerciements}
\addcontentsline{toc}{chapter}{Remerciements}
Texte de remerciements a personnaliser.

\chapter*{Resume}
\addcontentsline{toc}{chapter}{Resume}
Ce rapport presente la conception et la realisation d'une plateforme web de tourisme intelligent pour des hotels tunisiens. La solution combine une interface Next.js, un chatbot base sur IA generative avec approche RAG, un module analytics et un espace d'administration securise.

\chapter*{Abstract}
\addcontentsline{toc}{chapter}{Abstract}
This report presents the design and implementation of an intelligent tourism web platform for Tunisian hotels. The solution combines a Next.js interface, a generative AI chatbot with a RAG approach, an analytics module, and a secure administration space.

\chapter*{Liste des abreviations}
\addcontentsline{toc}{chapter}{Liste des abreviations}
\begin{itemize}[leftmargin=2cm]
  \item[IA] Intelligence Artificielle
  \item[LLM] Large Language Model
  \item[RAG] Retrieval-Augmented Generation
  \item[JWT] JSON Web Token
  \item[SSE] Server-Sent Events
  \item[KPI] Key Performance Indicator
  \item[API] Application Programming Interface
\end{itemize}

\tableofcontents
\listoffigures
\listoftables

\clearpage
\pagenumbering{arabic}

\chapter*{Introduction generale}
\addcontentsline{toc}{chapter}{Introduction generale}

Dans un contexte touristique de plus en plus numerise, les voyageurs attendent des reponses rapides, personnalisees et multilingues. Les hotels doivent ainsi offrir des experiences digitales modernes capables d'ameliorer la relation client avant, pendant et apres le sejour.

Notre projet de fin d'etudes s'inscrit dans cette dynamique. Il vise a concevoir une plateforme web intelligente permettant aux hotels tunisiens de centraliser leurs informations, d'assister les visiteurs via un chatbot base sur l'IA generative et de suivre les interactions utilisateurs a travers un module analytics.

La solution proposee repose sur une architecture moderne : Next.js pour l'application web, PostgreSQL pour la persistance des donnees, Redis pour le cache et les sessions, ainsi qu'une integration LLM pour la generation de reponses contextualisees via une approche RAG.

Afin d'atteindre ces objectifs, ce rapport est organise comme suit :
\begin{itemize}
  \item Le premier chapitre presente le contexte general du projet.
  \item Le deuxieme chapitre detaille l'analyse des besoins et l'architecture.
  \item Le troisieme chapitre expose la conception et la modelisation.
  \item Le quatrieme chapitre decrit l'implementation du chatbot intelligent.
  \item Le cinquieme chapitre traite le module analytics et l'espace administrateur.
  \item Le sixieme chapitre couvre la securite, les tests et le deploiement.
\end{itemize}

\chapter{Contexte general du projet}

\section{Introduction}
Ce chapitre presente le cadre general de notre projet, l'organisme d'accueil, l'analyse de l'existant, la problematique et la solution proposee.

\section{Presentation de l'organisme d'accueil}
\textit{[A completer avec les informations exactes de l'entreprise d'accueil : historique, missions, organisation, departement d'affectation.]}

\section{Contexte du domaine}
Le secteur hotelier tunisien connait une evolution importante en matiere de digitalisation. Les etablissements cherchent a ameliorer leur visibilite, leur qualite de service et leur capacite a repondre rapidement aux attentes des touristes nationaux et internationaux.

\section{Analyse de l'existant}
\subsection{Etude de l'existant}
Les approches traditionnelles reposent souvent sur des FAQ statiques, des supports disperses et une interaction limitee avec les utilisateurs.

\subsection{Problematique}
Les principales limites identifiees sont :
\begin{itemize}
  \item absence d'assistance conversationnelle contextuelle pour chaque hotel ;
  \item manque de personnalisation selon la langue et le profil client ;
  \item faible exploitation analytique des interactions utilisateurs ;
  \item difficultes de gestion centralisee des contenus hoteliers.
\end{itemize}

\subsection{Solution proposee}
Nous proposons une plateforme web intelligente qui integre un chatbot RAG multilingue, un tableau de bord analytics et un espace administrateur securise pour gerer les donnees hotelieres.

\section{Methodologie et planification}
\textit{[A completer : methode de gestion de projet choisie (Scrum/Kanban), planning, jalons, livrables.]}

\section{Conclusion}
Ce chapitre a pose le contexte du projet et justifie la solution retenue. Le chapitre suivant detaille les besoins fonctionnels et non fonctionnels ainsi que l'architecture cible.

\chapter{Analyse des besoins et architecture du systeme}

\section{Introduction}
Ce chapitre presente les besoins metiers, les exigences techniques et l'architecture globale de la solution.

\section{Identification des acteurs}
\begin{itemize}
  \item \textbf{Touriste / Client} : consulte les informations hotelieres et interagit avec le chatbot.
  \item \textbf{Administrateur d'hotel} : met a jour les parametres, evenements et contenus.
  \item \textbf{Super administrateur} : supervise l'acces, la securite et les indicateurs globaux.
\end{itemize}

\section{Besoins fonctionnels}
\begin{itemize}
  \item authentification et verification des acces administrateurs ;
  \item consultation et edition des informations par hotel ;
  \item chatbot IA avec reponses contextuelles et multilingues ;
  \item streaming de reponses et conservation de contexte conversationnel ;
  \item collecte et visualisation des statistiques d'utilisation.
\end{itemize}

\section{Besoins non fonctionnels}
\begin{itemize}
  \item securite (JWT, validation des entrees, sanitation, rate limiting) ;
  \item performance (cache Redis, optimisation des routes API) ;
  \item maintenabilite (architecture modulaire) ;
  \item evolutivite (ajout de nouveaux hotels et modules) ;
  \item disponibilite et robustesse en production.
\end{itemize}

\section{Architecture generale}
L'architecture est basee sur Next.js (App Router) avec des API Routes cote serveur. Les donnees sont persistees dans PostgreSQL, le cache et les sessions sont geres via Redis, et la couche IA est assuree via un modele LLM integre dans le pipeline de chat.

\begin{figure}[H]
  \centering
  \begin{tikzpicture}[
    node distance=1.2cm and 1.4cm,
    box/.style={rectangle,draw,rounded corners,align=center,minimum width=3.2cm,minimum height=0.9cm},
    db/.style={cylinder,draw,shape border rotate=90,aspect=0.25,align=center,minimum height=1.1cm,minimum width=2cm},
    >={Latex}
  ]
    \node[box] (client) {Client\\(Interface Web)};
    \node[box,below=of client] (next) {Next.js 14\\Frontend + API Routes};
    \node[box,below left=of next,xshift=-0.5cm] (llm) {Service IA\\Groq LLM};
    \node[box,below=of next] (rag) {Base de connaissance\\RAG};
    \node[db,below right=of next,xshift=0.5cm] (pg) {PostgreSQL};
    \node[db,right=of pg,xshift=0.6cm] (redis) {Redis};
    \node[box,right=of next,xshift=1.0cm] (admin) {Dashboard Admin\\Analytics};

    \draw[->] (client) -- (next);
    \draw[->] (next) -- (llm);
    \draw[->] (next) -- (rag);
    \draw[->] (next) -- (pg);
    \draw[->] (next) -- (redis);
    \draw[->] (admin) -- (next);
  \end{tikzpicture}
  \caption{Architecture generale de la plateforme}
\end{figure}

\section{Technologies adoptees}
\begin{table}[H]
  \centering
  \caption{Technologies principales du projet}
  \begin{tabular}{p{4cm}p{4cm}p{6cm}}
    \toprule
    \textbf{Couche} & \textbf{Technologie} & \textbf{Role} \\
    \midrule
    Frontend & Next.js 14 + React 18 & Interface utilisateur et navigation \\
    Backend & API Routes Next.js & Logique serveur et endpoints REST \\
    Base de donnees & PostgreSQL & Stockage persistant \\
    Cache & Redis & Cache applicatif et sessions \\
    IA & Groq + LLM & Generation de reponses \\
    Visualisation & Recharts & Dashboards analytics \\
    \bottomrule
  \end{tabular}
\end{table}

\section{Conclusion}
Apres l'etude des besoins et de l'architecture cible, nous passons dans le chapitre suivant a la conception detaillee et a la modelisation des donnees.

\chapter{Conception et modelisation}

\section{Introduction}
Ce chapitre decrit la conception de la solution, la structuration des donnees et les principaux flux fonctionnels.

\section{Conception fonctionnelle}
La solution s'articule autour de deux parcours principaux :
\begin{itemize}
  \item parcours client (consultation, profilage, chat intelligent) ;
  \item parcours administrateur (authentification, gestion contenus, suivi analytics).
\end{itemize}

\section{Modelisation de la base de donnees}
La base PostgreSQL est concue pour centraliser les informations hotelieres et les traces d'utilisation :
\begin{itemize}
  \item \texttt{hotels}, \texttt{hotel\_settings} ;
  \item \texttt{conversations}, \texttt{messages} ;
  \item \texttt{guest\_profiles}, \texttt{analytics\_events} ;
  \item \texttt{reservations}, \texttt{events}, \texttt{attractions} ;
  \item \texttt{admin\_users}.
\end{itemize}

\section{Conception des interfaces}
Les ecrans principaux incluent :
\begin{itemize}
  \item page d'accueil multi-hotels ;
  \item page detail hotel ;
  \item interface de chat temps reel ;
  \item tableau de bord administrateur ;
  \item interface analytics.
\end{itemize}

\section{Flux de donnees}
Le flux d'un message suit les etapes suivantes :
\begin{enumerate}
  \item reception et validation de la requete ;
  \item enrichment via contexte hotel + historique session ;
  \item generation de reponse par le modele IA ;
  \item persistence des traces et evenements analytics.
\end{enumerate}

\section{Conclusion}
Cette conception fournit une base solide pour la mise en oeuvre technique. Le chapitre suivant detaille l'implementation du chatbot intelligent et son integration RAG.

\chapter{Implementation du chatbot intelligent (RAG)}

\section{Introduction}
Ce chapitre presente l'implementation de la composante IA, coeur de la solution proposee.

\section{Architecture conversationnelle}
Le chatbot repose sur une approche RAG qui combine :
\begin{itemize}
  \item un modele de langage pour la generation ;
  \item une base de connaissances contextuelle par hotel ;
  \item une memoire de session pour la continuite des echanges.
\end{itemize}

\section{Pipeline de traitement d'un message}
\begin{enumerate}
  \item validation et nettoyage de l'entree utilisateur ;
  \item verification des limites d'usage (rate limiting) ;
  \item chargement du contexte hotelier et de l'historique ;
  \item construction du prompt ;
  \item appel LLM et retour de la reponse (JSON ou streaming SSE) ;
  \item enregistrement d'evenements analytiques.
\end{enumerate}

\section{Multilinguisme et personnalisation}
La plateforme prend en charge plusieurs langues afin d'ameliorer l'accessibilite internationale. Les reponses sont adaptees au contexte hotelier et au profil utilisateur.

\section{Exemple de logique de traitement}
\begin{lstlisting}[caption={Pseudo-flux de traitement d'une requete chat}]
1) Validate input
2) Apply rate-limit
3) Load hotel settings + session context
4) Build RAG prompt
5) Query model
6) Stream/return answer
7) Track analytics event
\end{lstlisting}

\begin{figure}[H]
  \centering
  \begin{tikzpicture}[
    node distance=1.0cm and 0.8cm,
    box/.style={rectangle,draw,rounded corners,align=center,minimum width=2.4cm,minimum height=0.8cm},
    >={Latex}
  ]
    \node[box] (u) {Utilisateur};
    \node[box,right=of u] (api) {API Chat};
    \node[box,right=of api] (val) {Validation\\+ Rate Limit};
    \node[box,right=of val] (ctx) {Contexte RAG\\(Hotel + Session)};
    \node[box,right=of ctx] (model) {LLM Groq};
    \node[box,below=of ctx] (track) {Tracking\\Analytics};

    \draw[->] (u) -- node[above]{message} (api);
    \draw[->] (api) -- (val);
    \draw[->] (val) -- (ctx);
    \draw[->] (ctx) -- (model);
    \draw[->] (model) -- node[above]{reponse} (api);
    \draw[->] (api) -- node[above]{SSE/JSON} (u);
    \draw[->] (api) -- (track);
  \end{tikzpicture}
  \caption{Diagramme de sequence simplifie du traitement d'un message chatbot}
\end{figure}

\section{Conclusion}
L'integration du chatbot RAG apporte une assistance intelligente et contextualisee. Le chapitre suivant aborde les modules analytics et administration.

\chapter{Module analytics et espace administrateur}

\section{Introduction}
Ce chapitre decrit la supervision fonctionnelle de la plateforme via les tableaux de bord et l'espace d'administration.

\section{Collecte des evenements analytics}
Le systeme enregistre les interactions importantes :
\begin{itemize}
  \item messages envoyes et reponses recues ;
  \item profils visiteurs (tranche d'age, nationalite, motif) ;
  \item actions d'administration.
\end{itemize}

\section{KPI et visualisations}
Les indicateurs permettent d'analyser :
\begin{itemize}
  \item usage du chatbot ;
  \item repartition demographique des utilisateurs ;
  \item activite par hotel et par periode ;
  \item performance des pages et interactions cles.
\end{itemize}

\section{Gestion administrative}
L'interface admin permet :
\begin{itemize}
  \item la gestion des parametres hoteliers ;
  \item la gestion des evenements et attractions ;
  \item le suivi des reservations ;
  \item la maintenance des contenus dynamiques.
\end{itemize}

\section{Authentification et controle d'acces}
L'acces administrateur est protege par JWT et verification serveur. Les routes sensibles sont controlees pour empecher les acces non autorises.

\section{Conclusion}
Ce module transforme les donnees d'usage en informations exploitables pour le pilotage operationnel. Le chapitre suivant traite des aspects securite, tests et deploiement.

\chapter{Securite, tests et deploiement}

\section{Introduction}
Ce chapitre presente les mecanismes de securite, la validation de la solution et son deploiement.

\section{Securite applicative}
Les mesures mises en place incluent :
\begin{itemize}
  \item hachage des mots de passe ;
  \item authentification JWT et verification de session ;
  \item validation des entrees (schema + sanitation) ;
  \item limitation du debit des requetes via Redis ;
  \item protection des routes administratives.
\end{itemize}

\section{Performance et scalabilite}
L'architecture exploite :
\begin{itemize}
  \item cache Redis pour reduire la latence ;
  \item appels non bloquants pour le tracking analytics ;
  \item streaming SSE pour ameliorer la reactivite percue ;
  \item approche modulaire favorisant la montee en charge.
\end{itemize}

\section{Strategie de tests}
La validation de la plateforme s'appuie sur :
\begin{itemize}
  \item tests des endpoints API ;
  \item tests fonctionnels par role ;
  \item tests de scenarios chat ;
  \item tests de non-regression sur les fonctionnalites critiques.
\end{itemize}

\section{Deploiement}
Le deploiement cible combine :
\begin{itemize}
  \item application Next.js en production ;
  \item PostgreSQL pour la persistence ;
  \item Redis pour cache et sessions ;
  \item variables d'environnement securisees.
\end{itemize}

\section{Conclusion}
Ces choix techniques garantissent une solution fiable, securisee et exploitable en contexte reel.

\chapter*{Conclusion generale et perspectives}
\addcontentsline{toc}{chapter}{Conclusion generale et perspectives}

Ce projet a permis de concevoir une plateforme touristique intelligente repondant aux besoins d'information, d'assistance et de pilotage des hotels tunisiens. Notre demarche a couvert l'ensemble du cycle, de l'analyse des besoins a la mise en production d'une solution web basee sur l'IA generative.

La valeur principale de la solution repose sur l'integration coherente de plusieurs composantes : chatbot contextualise via RAG, espace administrateur securise, et module analytics pour l'aide a la decision. Cette integration renforce a la fois l'experience utilisateur et l'efficacite operationnelle des etablissements.

En perspective, plusieurs evolutions peuvent etre envisagees :
\begin{itemize}
  \item extension mobile pour iOS et Android ;
  \item support vocal multilingue (speech-to-text / text-to-speech) ;
  \item enrichissement de la personnalisation basee sur l'historique utilisateur ;
  \item integration de nouveaux services (paiement, transport, activites partenaires).
\end{itemize}

En resume, ce travail confirme la pertinence d'une approche IA appliquee au tourisme, avec une architecture moderne, evolutive et orientee usage reel.


\bibliographystyle{plain}
\bibliography{references}

\end{document}
